\section*{Original geometry background needed for the construction (journal pp.23--25)}
\paragraph{Morphisms.}
Let $\mathcal V_1\subseteq\Aspace^n$ and $\mathcal V_2\subseteq\Aspace^m$ be affine varieties. A morphism from $\mathcal V_1$ to $\mathcal V_2$ is a map $f:\mathcal V_1\to\mathcal V_2$ that is a restriction of a polynomial map $\Aspace^n\to\Aspace^m$. Such a morphism $f$ is associated with a ring homomorphism $f^\sharp:\FF[\mathcal V_2]\to\FF[\mathcal V_1]$, making $\FF[\mathcal V_1]$ an algebra over $\FF[\mathcal V_2]$. We say $f$ is finite if $\FF[\mathcal V_1]$ is finitely generated as an $\FF[\mathcal V_2]$-module.
Let $f:\mathcal V_1\to\mathcal V_2$ be a map between projective varieties $\mathcal V_1$ and $\mathcal V_2$. We say $f$ is a morphism from $\mathcal V_1$ to $\mathcal V_2$ if there exists a collection of open subsets $\{U_i\}_{i\in I}$ of $\mathcal V_2$ such that $\mathcal V_2=\bigcup_{i\in I}U_i$ (i.e., $\{U_i\}_{i\in I}$ is an open cover of $\mathcal V_2$) and for each $i\in I$, the restriction $f|_{f^{-1}(U_i)}:f^{-1}(U_i)\to U_i$ is a morphism between affine varieties. Furthermore, if each $f|_{f^{-1}(U_i)}$ is finite, then we say $f$ is finite. Finiteness does not depend on the choice of the affine open cover. Namely, if $f:\mathcal V_1\to\mathcal V_2$ is a finite morphism between projective varieties $\mathcal V_1$ and $\mathcal V_2$, and $U$ is an open subset of $\mathcal V_2$ such that $f|_{f^{-1}(U)}:f^{-1}(U)\to U$ is a morphism between affine varieties, then $f|_{f^{-1}(U)}$ is also finite.
The image of a morphism $f:\mathcal V_1\to\mathcal V_2$ is denoted by $\operatorname{Im}(f)$ or $f(\mathcal V_1)$. The image of a closed set under a finite morphism is still closed. The composition of two finite morphisms is still finite.
\paragraph{Dimension.}
The dimension of an irreducible variety $\mathcal V$, denoted by $\dim(\mathcal V)$, is the largest integer $m$ such that there exists a chain of irreducible varieties $\emptyset\subsetneq\mathcal V_0\subsetneq\mathcal V_1\subsetneq\cdots\subsetneq\mathcal V_m=\mathcal V$. More generally, the dimension of a nonempty variety is the maximal dimension of its irreducible components. We define the dimension of an empty set to be $-\infty$. A variety is equidimensional if its irreducible components have the same dimension.
If $\pi:\mathcal V\to\mathcal V'$ is a finite morphism, then $\dim(\mathcal V)=\dim(\pi(\mathcal V))$.
\paragraph{Projective closure.}
The affine $n$-space $\Aspace^n$ may be regarded as an open subset of $\PP^n$ via the map $(x_1,\ldots,x_n)\mapsto(x_1,\ldots,x_n,1)$. The complement $H_\infty:=\PP^n\setminus\Aspace^n$ is a hyperplane of $\PP^n$ defined by $X_{n+1}=0$, called the hyperplane at infinity. For an affine subvariety $\mathcal V$ of $\Aspace^n\subseteq\PP^n$, the smallest projective subvariety of $\PP^n$ containing $\mathcal V$ is the projective closure of $\mathcal V$, which we denote by $\mathcal V^{\rm cl}$. It is known that $\mathcal V^{\rm cl}\cap\Aspace^n=\mathcal V$, $\dim(\mathcal V^{\rm cl})=\dim(\mathcal V)$, and $\deg(\mathcal V^{\rm cl})=\deg(\mathcal V)$.
\paragraph{Joins of disjoint projective varieties.}
For two distinct points $p,q\in\PP^n$, denote by $\overline{pq}$ the unique projective line passing through them. For two disjoint projective subvarieties $\mathcal V_1,\mathcal V_2\subseteq\PP^n$, define the join $J(\mathcal V_1,\mathcal V_2)$ of $\mathcal V_1$ and $\mathcal V_2$ as
\[J(\mathcal V_1,\mathcal V_2):=\bigcup_{p\in\mathcal V_1,q\in\mathcal V_2}\overline{pq}.\]
\begin{result}{Lemma 2.6 (Harris 1992, Examples 6.17, 11.36, and 18.17)}
$J(\mathcal V_1,\mathcal V_2)$ is a subvariety of $\PP^n$ of dimension $\dim(\mathcal V_1)+\dim(\mathcal V_2)+1$ and degree at most $\deg(\mathcal V_1)\cdot\deg(\mathcal V_2)$.
\end{result}
\begin{result}{Lemma 2.7 (Harris 1992, Exercise 11.6 and Corollary 18.5)}
Let $\mathcal V$ be a nonempty equidimensional subvariety of $\PP^n$ and $H$ a hypersurface of $\PP^n$ not containing an irreducible component of $\mathcal V$. Then $\mathcal V\cap H$ is an equidimensional subvariety of dimension $\dim(\mathcal V)-1$ and degree at most $\deg(\mathcal V)\cdot\deg(H)$ (or an empty set if $\dim(\mathcal V)=0$).
\end{result}
\begin{result}{Corollary 2.8}
Let $\mathcal V$ be a subvariety of $\PP^n$ of dimension $r$, where $0\le r<n$. Then there exists an $(n-r-1)$-subspace $W$ disjoint from $\mathcal V$.
\end{result}
\begin{proof}
It suffices to show that there exist hyperplanes $H_1,\ldots,H_{r+1}$, such that $\mathcal V_i:=\mathcal V\cap(\bigcap_{j=1}^iH_j)$ is empty for some $i\le r+1$. We may inductively choose each $H_i$ such that $C\not\subseteq H_i$ for every irreducible component $C$ of $\mathcal V_{i-1}$, so that $\dim(\mathcal V_i)\le\dim(\mathcal V_{i-1})-1$ by Lemma 2.7. So $\mathcal V_i=\emptyset$ for some $i\le r+1$.
\end{proof}
\begin{result}{Lemma 2.9 (Shafarevich 1994, Section I.6.2, Theorem 6)}
Suppose $\mathcal V_1$ and $\mathcal V_2$ are subvarieties of $\PP^n$ and $\dim(\mathcal V_1)+\dim(\mathcal V_1)\ge n$. Then $\mathcal V_1\cap\mathcal V_2\ne\emptyset$ and $\dim(\mathcal V_1\cap\mathcal V_2)\ge\dim(\mathcal V_1)+\dim(\mathcal V_1)-n$.
\end{result}
\section*{Original Section 3.1: Reducing to the Case of Equidimensional or Irreducible Varieties (journal pp.26--27)}
The following lemma states that to construct $k$-subspace families that are evasive for subvarieties of $\PP^n$, it suffices to consider equidimensional subvarieties of dimension $n-k-1$ (i.e., codimension $k+1$).
\begin{result}{Lemma 3.1}
Let $n,d\in\mathbb N^+$ and $k\in\{0,1,\ldots,n-1\}$. Let $\mathcal F$ be the family of all equidimensional subvarieties of $\PP^n$ of dimension $n-k-1$ and degree at most $d$. Then an $(\mathcal F,\epsilon)$-evasive $k$-subspace family is also $(n,d,\epsilon)$-evasive.
\end{result}
The proof of Lemma 3.1 is based on the following claim.
\begin{result}{Claim 3.2}
Let $\mathcal V$ be an irreducible subvariety of $\PP^n$. There exists a subvariety $\widetilde{\mathcal V}\subseteq\PP^n$ of dimension $n-k-1$ and degree at most $\deg(\mathcal V)$ such that any $k$-subspace of $\PP^n$ that evades $\widetilde{\mathcal V}$ also evades $\mathcal V$.
\end{result}
\begin{proof}
If $\dim(\mathcal V)=n-k-1$, then just let $\widetilde{\mathcal V}=\mathcal V$.
Now assume $\dim(\mathcal V)<n-k-1$. Let $t=(n-k-1)-\dim(\mathcal V)-1$ and let $\widetilde{\mathcal V}$ be the join of $\mathcal V$ and a $t$-subspace disjoint from $\mathcal V$ (which exists by Corollary 2.8). Then $\widetilde{\mathcal V}$ is a projective subvariety of dimension $n-k-1$ and degree at most $\deg(\mathcal V)$ by Lemma 2.6. Suppose $W$ is a $k$-subspace that evades $\widetilde{\mathcal V}$. Then $W$ is disjoint from $\widetilde{\mathcal V}\supseteq\mathcal V$. So $W$ also evades $\mathcal V$.
Finally, assume $\dim(\mathcal V)>n-k-1$. Let $t=\dim(\mathcal V)-(n-k-1)$. By Lemma 2.7, there exist $t$ hyperplanes $H_1,\ldots,H_t$ of $\PP^n$ such that $\mathcal V\cap\bigcap_{i=1}^tH_i$ is equidimensional of dimension $n-k-1$ and degree at most $\deg(\mathcal V)$. Let $\widetilde{\mathcal V}=\mathcal V\cap\bigcap_{i=1}^tH_i$. Suppose $W$ is a $k$-subspace that evades $\widetilde{\mathcal V}$. Then $W\cap\widetilde{\mathcal V}=(W\cap\mathcal V)\cap\bigcap_{i=1}^tH_i=\emptyset$. Again by Lemma 2.7, we have $\dim(W\cap\mathcal V)\le t-1=\dim(\mathcal V)+\dim(W)-n$. So $W$ also evades $\mathcal V$.
\end{proof}
\begin{proof}[Proof of Lemma 3.1]
Consider a projective subvariety $\mathcal V\subseteq\PP^n$ of degree at most $d$. Let $\mathcal V_1,\ldots,\mathcal V_s$ be the irreducible components of $\mathcal V$. For each $i\in[s]$, use Claim 3.2 to choose a projective subvariety $\widetilde{\mathcal V}_i\subseteq\PP^n$ of dimension $n-k-1$ and degree at most $\deg(\mathcal V_i)$ such that any $k$-subspace that evades $\widetilde{\mathcal V}_i$ also evades $\mathcal V_i$. Let $\widetilde{\mathcal V}=\bigcup_{i=1}^s\widetilde{\mathcal V}_i$. Then $\widetilde{\mathcal V}\in\mathcal F$. By construction, any $k$-subspace that evades $\widetilde{\mathcal V}$ also evades $\mathcal V$. It follows that an $(\mathcal F,\epsilon)$-evasive $k$-subspace family is also $(n,d,\epsilon)$-evasive.
\end{proof}
We further reduce to the case of irreducible varieties at the cost of blowing up the parameter $\epsilon$ by a factor of $d$. This is useful as we need irreducibility later in Lemma 3.13.
\begin{result}{Lemma 3.3}
Let $n,d\in\mathbb N^+$ and $k\in\{0,1,\ldots,n-1\}$. Let $\mathcal F'$ be the family of all irreducible subvarieties of $\PP^n$ of dimension $n-k-1$ and degree at most $d$. Then an $(\mathcal F',\epsilon)$-evasive $k$-subspace family is also an $(n,d,d\epsilon)$-evasive $k$-subspace family.
\end{result}
\begin{proof}
Let $\mathcal F$ be as in Lemma 3.1. Each $\mathcal V\in\mathcal F$ has at most $d$ irreducible components, which are all in $\mathcal F'$ since their degrees are bounded by $d$. By definition and the union bound, if a $k$-subspace family $\mathcal H$ is $(\mathcal F',\epsilon)$-evasive, then it is also $(\mathcal F,d\epsilon)$-evasive. Combining this with Lemma 3.1 proves the lemma.
\end{proof}
\section*{Original Chow forms in Stiefel coordinates (journal pp.31--32)}
We may also express the Chow form in Stiefel coordinates, i.e., in the entries of a generating matrix of a linear subspace. This expression has the advantage that it is an actual polynomial rather than a member of the abstract vector space $(R/I)_d$.
Formally, let $A^*$ be a $(k+1)\times(n+1)$ variable matrix whose $(i,j)$-th entry is a variable $Y_{i,j}$. Define the ring homomorphism
\[\phi^\sharp:R=\FF[X_S:S\in\binom{[n+1]}{k+1}]\to\FF[Y_{i,j}:i\in[k+1],j\in[n+1]]\]
that sends each variable $X_S$ to $\det(A^*_{[k+1],S})$. Define the Chow form of $\mathcal V$ in Stiefel coordinates to be
\[\widetilde R_{\mathcal V}:=\phi^\sharp(P_{\mathcal V})\in\FF[Y_{i,j}:i\in[k+1],j\in[n+1]]\]
where $P_{\mathcal V}\in R_d$ is a lift of $R_{\mathcal V}\in(R/I)_d$. Note that $I$ is precisely the kernel of $\phi^\sharp$. So $\widetilde R_{\mathcal V}$ is uniquely determined by $\mathcal V$ up to scalars. By construction, for any $W\in\GG(k+1,n+1)$ and generating matrix $A=(a_{i,j})_{i\in[k+1],j\in[n+1]}$ of $W$, we have $P_{\mathcal V}(\phi(W))=\widetilde R_{\mathcal V}(A):=\widetilde R_{\mathcal V}(a_{1,1},\ldots,a_{k+1,n+1})$. So $\widetilde R_{\mathcal V}$ vanishes at $A$ iff $\PP W\in\GG(k,n)$ intersects $\mathcal V$.
\paragraph{Chow forms in dual Stiefel coordinates.}
Similarly, we may express the Chow form in dual Stiefel coordinates, i.e., in the entries of a parity check matrix of a linear subspace.
More specifically, choose a homogeneous polynomial $Q_{\mathcal V}\in\FF[X_S:S\in\binom{[n+1]}{n-k}]$ that defines the set of $k$-subspaces intersecting $\mathcal V$ in terms of dual Pl\"ucker coordinates. As primal and dual Pl\"ucker coordinates are equivalent, $Q_{\mathcal V}$ can be obtained from the polynomial $P_{\mathcal V}$ above by simply negating and renaming variables. Next, compose $Q_{\mathcal V}$ with a ring homomorphism that substitutes dual Pl\"ucker coordinates with dual Stiefel coordinates. The resulting polynomial, which we denote by $\widetilde R^\vee_{\mathcal V}\in\FF[Y_{i,j}:i\in[n-k],j\in[n+1]]$, is called the Chow form of $\mathcal V$ in dual Stiefel coordinates.
We note that the Chow form $\widetilde R_{\mathcal V}$ in primal Stiefel coordinates is a homogeneous polynomial of degree $(k+1)d$ in $(k+1)(n+1)$ variables, whereas the Chow form $\widetilde R^\vee_{\mathcal V}$ in dual Stiefel coordinates is a homogeneous polynomial of degree $(n-k)d$ in $(n-k)(n+1)$ variables. This suggests that it is more convenient to use the Chow form in primal (resp. dual) Stiefel coordinates when $k$ is small (resp. $n-k$ is small).
